% !PRINTED none
% The source page between printed pp. 521 and 523 is unnumbered; its complete
% section-opening content is transcribed below.
\begin{center}
\gesperrt{F\"unfzehnter Abschnitt.}

\gesperrt{Cyklische Gleichungen.}

\rule{3cm}{0.4pt}

\S. 159.

\gesperrt{Cubische Gleichungen.}
\end{center}

Wir wollen von dem jetzt gewonnenen Standpunkte aus zun\"achst die
Aufl\"osung der Gleichungen dritten und vierten Grades
betrachten.

Die Gruppe der Permutationen von drei Ziffern besteht aus
sechs Elementen, n\"amlich aus der identischen Permutation 1, aus
zwei dreigliedrigen Cyklen und drei Transpositionen:
\weq{1}{\begin{array}{ccc} 1 & (0,\,1,\,2) & (0,\,2,\,1) \\ (1,\,2) & (2,\,0) & (0,\,1). \end{array}}
Die drei ersten
\[
1, \; (0,\,1,\,2) \quad (0,\,2,\,1)
\]
bilden die alternirende Gruppe, sie besteht aus den Potenzen der
cyklischen Permutation $\pi = (0, 1, 2)$:
\weq{2}{\pi^0 = 1, \; \pi, \; \pi^2}
und wird daher \gesperrt{cyklische Gruppe} genannt.

Es seien nun $\alpha, \alpha_1, \alpha_2$ die drei Wurzeln der cubischen
Gleichung:
\weq{3}{f(x) = x^3 - ax^2 + bx - c = 0.}
Legen wir den K\"orper $\Omega$ zu Grunde, der aus allen rationalen
 Functionen der unabh\"angigen Ver\"anderlichen $a, b, c$ besteht,
so ist (1) die Gruppe dieser Gleichung; wenn wir aber
\weq{4}{\sqrt{D} = (\alpha - \alpha_1)(\alpha - \alpha_2)(\alpha_1 - \alpha_2)}
%CONT>>
